\documentclass[12pt]{article}
\usepackage[utf8]{inputenc}
\usepackage[T1]{fontenc}
\usepackage{hyperref}
\title{Tanosy Language Resources for K-12 Curriculum Development}
\author{Compiled by Claude Sonnet 5 (claude-sonnet-5) for Andr\'as Kornai}
\date{}
\begin{document}
\maketitle

\section{Introduction}
Tanosy (Antanosy) is a variety of Malagasy (ISO 639-3 txy) spoken in southeastern Madagascar. The searches found no resources written specifically for Tanosy. Everything below concerns Malagasy in general. Whether it applies to Tanosy is not established.

\section{Primary Educational Resources}
No Tanosy primers, textbooks or school materials turned up in the searches.

\section{Linguistic and Reference Materials}
\begin{itemize}
\item Malagasy online dictionary (multiple dictionaries, with some glosses labelled ``[Tanosy]''): \url{https://malagasyword.org/bins/teny2/dictionary}
\item Lexilogos collection of links to historical Malagasy dictionaries (Richardson, Freeman, Malzac, Flacourt): \url{https://www.lexilogos.com/english/malagasy_dictionary.htm}
\item ``The Malagasy language in the digital age'' (HAL), on language technology for Malagasy: \url{https://hal.archives-ouvertes.fr/hal-01237509}
\end{itemize}

\section{Religious and Community Resources}
Scripture Earth lists Malagasy Bible editions under the Plateau Malagasy code (plt), not Tanosy: \url{https://www.scriptureearth.org/00eng.php?iso=plt}

\section{Digital and Audio Resources}
No Tanosy audio or video resources were found.

\section{Curriculum Development Context}
Standard Malagasy materials are the only starting point found. A Tanosy curriculum would need new, locally gathered material, and further searches such as ``Antanosy dialect'' are advisable.
\end{document}